\documentclass[10pt,a4paper]{amsart}
\usepackage[margin=19mm]{geometry}
\usepackage{amsmath,amssymb,mathtools}
\usepackage[scr=rsfs,cal=euler]{mathalfa}
\usepackage{xfrac}
\usepackage[unicode,hidelinks,bookmarks=false]{hyperref}
\newcommand{\ie}{i.e.\ }
\newcommand{\cf}{cf.\ }
\newcommand{\resp}{respectively\ }
\newcommand{\Subsection}{\subsection}
\providecommand{\nobreakdash}{\nobreak\hbox{-}}
\setlength{\emergencystretch}{2em}
\multlinegap0pt
\usepackage{bm}
\usepackage{bbm}
\usepackage{dsfont}
\usepackage{xspace}
\usepackage{cancel}
\usepackage{mathtools}
\usepackage{amsthm}
\makeatletter
\@ifundefined{theorem}{\newtheorem{theorem}{Theorem}}{}
\@ifundefined{claim}{\newtheorem{claim}[theorem]{Claim}}{}
\makeatother
\newcommand{\IF}{\mathbb{F}}
\newcommand{\IN}{\mathbb{N}}
\newcommand{\SL}{\textsf{\textup{SL}}}
\newcommand{\Sym}{\textup{Sym}}
\newcommand{\oIF}{{\overline{\mathbb{F}}}}
\renewcommand{\Im}{\mathsf{Im}}
\title[Field-independent Kronecker-plethysm isomorphisms]{Field-independent Kronecker-plethysm isomorphisms}
\author{Christian Ikenmeyer, Heidi Omar, Dimitrios Tsintsilidas}
\date{Source: arXiv:2509.10069v2 (CC BY 4.0)}
\begin{document}
\maketitle

\noindent\emph{Source and licence.} Field-independent Kronecker-plethysm isomorphisms. Version arXiv:2509.10069v2 (CC BY 4.0), distributed under the Creative Commons Attribution 4.0 International licence, as recorded in the arXiv licence metadata for this exact version and shown on the abstract page https://arxiv.org/abs/2509.10069v2. Full rights notice: licenses/arxiv-2509.10069-CC-BY-4.0.txt.\par\smallskip

\noindent\textit{Editorial scope.} Counted unit: the Claim whose source label is \texttt{cla:imageelemsym} in the article (source0.tex lines 564--567, source label \texttt{cla:imageelemsym}), delivered together with the article's own complete original proof of it (source0.tex lines 568--618, one \texttt{proof} environment of 51 lines). No mathematics is written, repaired or re-derived: statement and proof are the article's own text line for line. Required context retained uncounted: eq:PhiMe (lines 462--464); cla:multipleofell (lines 513--515); cla:symmetricinmu (lines 546--548); eq:degmui (lines 560--562). Every definition, display and label of those lines is retained unchanged. Omitted from this package and not redistributed: 1--461, 465--512, 516--545, 549--559, 563--563, 619--1358. Nothing inside a retained excerpt is omitted or altered: every retained body is delivered exactly as the article's lines are written, so no omission receipt of any kind accompanies this package. Citations: the retained excerpts carry no citation command at all, so no bibliography entry of the article is transcribed and no reference is redistributed. Reference adaptation: none was needed and none was made; every reference inside the delivered unit resolves inside it, so no unresolved reference or citation is produced. Printed slips kept literal: no printed word, symbol, hypothesis or number of the counted statement or of its proof was altered. Layout and class: the article's own class is \texttt{extarticle} with packages including amsmath, amssymb, amsthm, mathrsfs, mathtools, geometry, comment, url, verbatim, ytableau, youngtab, tikz, float, hyperref; the wrapper is the lane's pinned \texttt{amsart} editorial wrapper at 10pt on A4, which restates the theorem-like environments the retained text needs and adds the article's own macro definitions that the retained text needs (\texttt{\textbackslash IF}, \texttt{\textbackslash IN}, \texttt{\textbackslash Im}, \texttt{\textbackslash SL}, \texttt{\textbackslash Sym}, \texttt{\textbackslash oIF}), with extra wrapper packages (bm, bbm, dsfont, xspace, cancel, mathtools). The delivered numbering is the unit's own. Assets: no retained span contains a figure environment, a graphics-inclusion command, a PDF-page inclusion, a table or a TikZ picture; no image file is copied into this package, the package declares no asset dependency and no image file is redistributed with it. Licence: version arXiv:2509.10069v2 of this article is distributed under the Creative Commons Attribution 4.0 International licence, as recorded in the arXiv licence metadata for this exact version and shown on the abstract page https://arxiv.org/abs/2509.10069v2. A full rights notice is delivered at licenses/arxiv-2509.10069-CC-BY-4.0.txt. Characters: every retained excerpt is pure ASCII, so the delivered engine, XeLaTeX with its default Latin Modern text font, reports no missing character.
\begin{equation}\label{eq:PhiMe}
\Phi(M_\ell(k)) = \nu^{k} e_{\ell-k}.
\end{equation}

\begin{claim}\label{cla:multipleofell}
For every $p \in (\Sym^{\bullet}(\IF^{\ell\times\ell\times 2}))^{\SL_\ell(\oIF)\times \SL_\ell(\oIF)}$, we have that $d:=\deg(p)$ is a multiple of $\ell$.
\end{claim}

\begin{claim}\label{cla:symmetricinmu}
$\Phi(p)$ is symmetric in the variables $\mu_1,\ldots,\mu_\ell$.
\end{claim}

\begin{equation}\label{eq:degmui}
\textstyle\deg_{\mu_i}(p)\leq\frac{d}{\ell}.
\end{equation}

\begin{claim}
\label{cla:imageelemsym}
$\Im(\Phi) = \IF\left[\nu^k e_{\ell-k}\mid 0\leq k \leq \ell\right]$.
\end{claim}

\begin{proof}
We have
$\Phi(M_\ell(k)) = \nu^k e_{\ell-k}$, see \eqref{eq:PhiMe}, so it suffices to show that
$\Im(\Phi) \subseteq \IF\left[\nu^k e_{\ell-k}\mid 0\leq k \leq \ell\right]$.
Since the group action of $\SL_\ell(\oIF)\times \SL_\ell(\oIF)$ preserves degrees, every invariant decomposes into a sum of homogeneous invariants.
Let $p \in (\Sym^{\bullet}(\IF^{\ell\times\ell\times 2}))^{\SL_\ell(\oIF)\times \SL_\ell(\oIF)}$
be homogeneous of some degree $d$.
It remains to show that $\Phi(p) \in \IF\left[\nu^k e_{\ell-k}(\mu_1,\dots,\mu_\ell)\mid 0\leq k \leq \ell\right]$.
We collect powers of $\nu$ in $\Phi(p)$ and express the symmetric part (Claim~\ref{cla:symmetricinmu}) as a polynomial in elementary symmetric polynomials:
\begin{equation}
\label{eq:pprimeine}
\Phi(p)(\nu,\mu_1,\dots,\mu_\ell)
 \ = \ 
\sum_{\delta=0}^{d} \nu^\delta \sum_{\mathbf{d}=(d_1,\ldots,d_\ell)} \alpha_{\mathbf{d}} \, e_1^{d_1} \cdots e_\ell^{d_\ell}
\end{equation}
for constants $\alpha_{\mathbf{d}}\in \IF$,
and the sum for every $\delta$ is over
$(d_1,\ldots,d_\ell) \in \IN_0^\ell$ with
\begin{equation}
\label{eq:1}
\textstyle
\sum_{k=1}^\ell d_k \cdot k = d-\delta,
\end{equation}
since $\deg (e_k)=k$.
Combining \eqref{eq:degmui} with the fact that $\deg_{\mu_i}(e_k) = 1$, we also obtain
\begin{equation}
\label{eq:2}
\textstyle
\sum_{k=1}^\ell d_k \leq \frac{d}{\ell}.
\end{equation}
Therefore, for every $\delta$ we have
\[\textstyle
\delta
\stackrel{\eqref{eq:1}}{=}
d - \sum_{k=1}^\ell d_k \cdot k
\stackrel{\eqref{eq:2}}{\geq}
\sum_{k=1}^\ell d_k (\ell-k),
\]
hence $\delta' := \delta-\sum_{k=1}^\ell d_k (\ell-k) \geq 0$.
Therefore we can rewrite the monomial $\nu^\delta e_1^{d_1}e_2^{d_2}\cdots e_\ell^{d_\ell}$
from \eqref{eq:pprimeine}
as
\begin{equation}
\label{eq:newmonomial}
\nu^{\delta'} (\nu^{\ell-1}e_1)^{d_1}(\nu^{\ell-2}e_2)^{d_2}\cdots (\nu e_{\ell-1})^{d_{\ell-1}} (\nu^0 e_\ell)^{d_\ell}.
\end{equation}
Since $d$ is a multiple of $\ell$ (Claim~\ref{cla:multipleofell}),
we can also see from \eqref{eq:1} that $\delta'= d - \ell\cdot \sum_{k=1}^\ell d_k$, so $\delta'$ is a multiple of $\ell$,
hence $\nu^{\delta'} = (\nu^\ell e_0)^{d_0}$ with $d_0 = \frac{d}{\ell}-\sum_{k=1}^\ell d_k$.
Overall, the monomial in \eqref{eq:newmonomial} is contained in $\IF\left[\nu^k e_{\ell-k}(\mu_1,\dots,\mu_\ell)\mid 0\leq k \leq \ell\right]$, which implies that $\Phi(p)$ is also contained there, as desired.
\end{proof}

\end{document}
